%% =====================================================================
%%  T-20-01  Measured Pressure
%%  -------------------------------------------------------------------
%%  One idea per file. Core is mandatory. Slots in canonical order.
%%  Max 700 words. No layout commands. No filler. New source material
%%  for this entry goes in sources/<BOOK>/T-20-01.tex -- never by
%%  rewriting this file (docs/RULES.md R0.1).
%%  =====================================================================
%% @kind: topic
%% @id: T-20-01
%% @title: Measured Pressure
%% @subtitle: A little pressure has its place --- but seldom, and once
%% @chapter: c20
%% @order: 10
%% @status: stable
%% @sources: B1
%% @updated: 2026-09-19
%% @allow-rewrite: B1 ingest 2026-09-19 -- committed scaffold replaced by the written entry (planned -> stable lifecycle)

\topic{T-20-01}{Measured Pressure}{A little pressure has its place --- but seldom, and once}

\begin{core}
Pressure --- a deadline, a stake, a single hard choice delivered flatly --- belongs in this craft the way a hammer belongs in a toolbox: real, necessary, and wrong for almost every job. It is reserved for the counterpart who is genuinely immovable and genuinely important, it is applied once, in full, at the chosen moment, and it ends the persuasion phase permanently.
\end{core}
\begin{mechanism}
The source's chapter is framed as the exception that proves the whole method: most people cannot be swayed by even masterly persuasion --- minds too fixed, factors beyond your control --- and wisdom is recognising those dead ends in advance so effort is not poured into unwinnable cases. For the rest, the instrument of last resort is the fell stroke: pressure applied when the target shows the yes-no state --- visibly caught between the comfort of agreeing and the fear of refusing. Its power is cumulative context, not volume: every prior hour of rapport, patience and credible reason is what makes one hard moment land as decision rather than as attack. Pressure delivered without that context is just coercion, and it produces the two failure products of coercion --- resentment, or the compliance that dies the instant supervision relaxes. It also cannot be repeated: the first stroke spends the surprise on which the second would depend.
\end{mechanism}
\begin{conditions}
Legitimate and effective only where all of these hold: the relationship justifies the damage; the matter is vital; persuasion has genuinely run its course; and you are willing to have the counterpart remember it forever. Everywhere else it is the amateur's tell --- pressure as a substitute for preparation is the most legible failure in the whole book.
\end{conditions}
\begin{application}
  \item Confirm the dead end honestly: three genuine persuasion attempts, differently built, or no stroke at all.
  \item Pick the moment of maximum yes-no tension --- not your peak anger, their peak ambivalence.
  \item Name the choice once, flatly, with its real consequence attached. No theatre, no repetition.
  \item Stop talking immediately after. Every word past the stroke discounts it.
  \item Have the follow-through ready. A stroke you do not carry out teaches the target that your words are theatre.
\end{application}
\begin{feedback}
  \item Landing: a long stillness, then a decision --- either answer. The stroke did its work as information.
  \item Landing: the relationship is colder but the terms are honest. That is the trade working as designed.
  \item Failing: they agree to end the conversation. You bought compliance and taught them your limit.
  \item Failing: they call it early --- \emph{is this a threat?} --- and the frame collapses into theatre.
\end{feedback}
\begin{failure}
The tool's dangers are structural. Escalation is one-way: a stroke that fails demands a bigger one, and the ladder ends somewhere illegal. The relationship survives only as a renegotiated arrangement --- fear supervises nothing once fear departs. And the practitioner pays too: the habit of reaching for pressure atrophies every instrument that made it unnecessary.
\end{failure}
\begin{countermeasures}
  \item Learn the amateur's signature: stakes named before rapport earned, deadlines with no origin, consequences narrated with relish. Walk.
  \item When pressed, slow the clock and get the pressure into writing --- coercion hates a record (T-24-03).
  \item Separate the real stake from the staged one: pressure is credible exactly once per relationship. Demand evidence for the second.
  \item In your own hands, audit the stroke you are proud of: would it survive being described aloud to the person's family?
\end{countermeasures}

\sources{B1}
\seesources
\seealso{T-20-02, T-24-03, T-16-02}
